\documentclass[11pt]{article}

\usepackage[margin=1in]{geometry}
\usepackage{amsmath, amssymb}
\usepackage{booktabs}
\usepackage{graphicx}
\usepackage{float}
\usepackage{hyperref}
\usepackage{authblk}
\usepackage{caption}

\title{Deep Learning Statistical Arbitrage:\\ An Empirical Study of Neural Trading Policies on Equity Residuals}
\author[1]{Hadis Surya Al Amin}
\affil[1]{\small Project report -- based on the framework of Guijarro-Ordonez, Pelger, and Zanotti (2021)}
\date{\today}

\begin{document}
\maketitle

\begin{abstract}
We study data-driven statistical arbitrage in U.S.\ equities by learning trading
policies directly from \emph{residual} return series. Residuals are extracted with
three families of factor models---IPCA, PCA, and the Fama--French five-factor
model---and a neural trading policy is trained end-to-end to maximize an
out-of-sample Sharpe ratio. We evaluate a convolution--transformer architecture
with a graph component (\texttt{CNNTransformer}) against feed-forward baselines
(\texttt{RawFFN}, \texttt{FourierFFN}, \texttt{OUFFN}). This report describes the
data, methodology, and empirical results, including risk-adjusted performance and
turnover across factor models.
\end{abstract}

\section{Introduction}
Classical statistical arbitrage exploits temporary mispricings in the residual
component of asset returns---the part left after systematic (factor) exposure is
removed. Rather than assuming a parametric mean-reverting process (e.g.\ an
Ornstein--Uhlenbeck model) and hand-tuning entry/exit thresholds, we follow the
\emph{Deep Learning Statistical Arbitrage} framework and learn the trading policy
directly from data. The policy maps a window of recent residual behavior to a
portfolio weight, and is optimized to maximize a risk-adjusted objective on
out-of-sample (OOS) data via rolling retraining.

\section{Data and Residual Construction}
The investable universe is defined by a market-capitalization cap proportion
($\texttt{cap\_proportion} = 0.01$). For each asset we form a residual time series
by removing common factors using one of three factor models:
\begin{itemize}
  \item \textbf{IPCA} --- Instrumented Principal Component Analysis (here $5$ factors),
  \item \textbf{PCA} --- rolling-window principal component analysis ($5$ factors),
  \item \textbf{Fama--French} --- the daily five-factor model.
\end{itemize}
Daily trading dates are aligned to the Fama--French research data over the OOS
period (1998--2017). When \texttt{use\_residual\_weights} is enabled, the residual
composition matrix is used to translate residual-space positions into asset-space
weights for realistic turnover and short-proportion accounting.

\section{Methodology}

\subsection{Trading policy models}
We compare four trading-policy architectures, all consuming a lookback window of
preprocessed residuals (cumulative-sum normalized):
\begin{itemize}
  \item \texttt{RawFFN} --- feed-forward network on raw residual windows (baseline).
  \item \texttt{FourierFFN} --- feed-forward network on Fourier features of the window.
  \item \texttt{OUFFN} --- feed-forward network with an Ornstein--Uhlenbeck-inspired layer.
  \item \texttt{CNNTransformer} --- a one-dimensional convolutional encoder followed by a
        multi-head self-attention (transformer) block, optionally augmented by a
        fully connected graph neural network (GNN) component to share information
        across the cross-section.
\end{itemize}
The reported \texttt{CNNTransformer} configuration uses a lookback of $30$ days,
dropout $0.25$, convolutional filters $[1,8]$ of size $2$, $4$ attention heads,
and a single GNN layer with hidden dimension $8$.

\subsection{Objective and training protocol}
Let $r_t$ denote the daily portfolio return generated by the policy. The policy
parameters $\theta$ are trained to maximize the in-sample Sharpe ratio
\begin{equation}
  \mathcal{L}(\theta) = \frac{\mathbb{E}[r_t]}{\sqrt{\operatorname{Var}[r_t]}},
\end{equation}
optionally net of transaction costs (\texttt{trans\_cost}) and holding/short costs
(\texttt{hold\_cost}). Training uses Adam (learning rate $10^{-3}$, batch size $32$)
for up to $100$ epochs. We employ a \textbf{rolling-window} protocol: the model is
retrained on a window of \texttt{length\_training}\,$=1000$ trading days and applied
to the subsequent block, with retraining every \texttt{retrain\_freq}\,$=125$ days,
so all reported performance is strictly out-of-sample.

\subsection{Evaluation metrics}
For each (model, factor model) pair we report the annualized Sharpe ratio
(daily Sharpe $\times \sqrt{252}$), mean daily return, return standard deviation,
turnover, and short proportion.

\section{Results}
Table~\ref{tab:summary} summarizes risk-adjusted performance across factor models,
sorted by annualized Sharpe ratio. Figure~\ref{fig:equity} shows the corresponding
out-of-sample equity curves (cumulative return, log scale).

\begin{table}[H]
  \centering
  \caption{Out-of-sample performance by trading policy and factor model.}
  \label{tab:summary}
  \small
  % Generated by: python3 tools/analyze_results.py
  % Produces ../results/summary_table.tex
  \IfFileExists{../results/summary_table.tex}{%
    \input{../results/summary_table.tex}%
  }{%
    \textit{Run \texttt{python3 tools/analyze\_results.py} to generate this table.}%
  }
\end{table}

\begin{figure}[H]
  \centering
  \IfFileExists{../results/equity_curves.png}{%
    \includegraphics[width=0.95\textwidth]{../results/equity_curves.png}%
  }{%
    \fbox{\parbox{0.6\textwidth}{\centering\vspace{2cm}
    \textit{equity\_curves.png will appear here after running the analysis script.}
    \vspace{2cm}}}%
  }
  \caption{Out-of-sample cumulative returns (equity curves) for each trading
  policy, log scale, normalized to a starting value of $1.0$.}
  \label{fig:equity}
\end{figure}

\subsection{Discussion}
% Fill in once results are downloaded and the table/figure are generated.
We discuss which factor-model residuals yield the strongest risk-adjusted returns,
the trade-off between Sharpe ratio and turnover, and the marginal value of the
transformer and GNN components relative to the feed-forward baselines.

\section{Conclusion}
We presented an end-to-end deep learning approach to statistical arbitrage on
equity residuals and benchmarked several neural trading-policy architectures across
IPCA, PCA, and Fama--French factor models under a rolling out-of-sample protocol.

\section*{Reproducibility}
Results are produced with \texttt{run\_train\_test.py} using the configuration files
in \texttt{configs/}. To regenerate the table and figure in this report:
\begin{verbatim}
python3 tools/analyze_results.py            # all models
python3 tools/analyze_results.py -m CNNTransformer
\end{verbatim}
To match the published paper, set \texttt{use\_residual\_weights: True} in the config.

\begin{thebibliography}{9}
\bibitem{dlsa} J.~Guijarro-Ordonez, M.~Pelger, and G.~Zanotti.
  \emph{Deep Learning Statistical Arbitrage}, 2021.
  \url{https://arxiv.org/abs/2106.04028}.
\end{thebibliography}

\end{document}
